\section{CMOS Logic}\label{cmos-logic}\index{cmos logic@CMOS Logic}

\textbf{Keywords:} CMOS Logic, Inverter, NAND, Static Logic, SR-Latch,
D-Latch, Flip-Flop, Tristate, AOI

\section{Analog transistor to digital
transistor}\label{analog-transistor-to-digital-transistor}

NMOS current (W = 0.4u L=0.15u) as a function of \[V_{GS}\] and
\[V_{DS}\]

\_\_dicex/lectures/l13/mos.py


{
\centering
\aicfig{media/transistor_log.png}
}

\emph{Figure 1: NMOS drain current on a log10 scale as a surface over
gate-source and drain-source voltage, for W = 0.4u and L = 0.15u}


{
\centering
\aicfig{media/transistor_lin.png}
}

\emph{Figure 2: The same NMOS drain current on a linear scale in mA over
gate-source and drain-source voltage}


{
\centering
\aicfig{media/analog.png}
}

\emph{Figure 3: The analog view of the NMOS: current equations for the
linear and saturation regions across weak, moderate and strong inversion
and mobility degradation}


{
\centering
\aicfig{media/digital.png}
}

\emph{Figure 4: The digital view of the same plane, where weak inversion
is treated as OFF and strong inversion as ON, and everything else is
ignored}

\begin{longtable}[]{@{}ccc@{}}
\toprule\noalign{}
Gate & NMOS & PMOS \\
\midrule\noalign{}
\endhead
\bottomrule\noalign{}
\endlastfoot
VDD & ON & OFF \\
VDD -\textgreater{} VSS & X & X \\
VSS -\textgreater{} VDD & X & X \\
VSS & OFF & ON \\
\end{longtable}

\begin{longtable}[]{@{}ccc@{}}
\toprule\noalign{}
Gate & NMOS & PMOS \\
\midrule\noalign{}
\endhead
\bottomrule\noalign{}
\endlastfoot
1 & ON & OFF \\
1 -\textgreater{} 0 & X & X \\
0 -\textgreater{} 1 & X & X \\
0 & OFF & ON \\
\end{longtable}

\section{CMOS static logic
assumptions}\label{cmos-static-logic-assumptions}

NMOS source is connected to low potential

\[ V_{GS} > V_{TH}\] when \[V_G = V_{DD}\]

PMOS source is connected to high potential

\[ V_{GS} < V_{TH}\] when \[V_G = 0\]


{
\centering
\aicfig{media/nand_tr_tikz.pdf}
}

\emph{Figure 5: Transistor schematic of a two input NAND, with PMOS A
and B in parallel to the supply and NMOS A and B in series to ground}


{
\centering
\aicfig{media/rules.pdf}
}

\emph{Figure 6: Two ways to wire a NOR - the accepted one with PMOS
pulling up and NMOS pulling down, and the rejected one with the device
types swapped so the sources sit at the wrong rail}

\section{Don't break rules unless you know exactly why it will be
OK}\label{dont-break-rules-unless-you-know-exactly-why-it-will-be-ok}

\section{Logic cells}\label{logic-cells}\index{logic cells@Logic cells}


{
\centering
\aicfig{media/binary_tikz.pdf}
}

\emph{Figure 7: Truth tables for NAND, NOR, AND and OR with inverted
inputs, illustrating the two De Morgan identities}

\subsection{CMOS static logic is
inverting}\label{cmos-static-logic-is-inverting}

\begin{longtable}[]{@{}cc@{}}
\toprule\noalign{}
A & Y \\
\midrule\noalign{}
\endhead
\bottomrule\noalign{}
\endlastfoot
1 & 0 \\
0 & 1 \\
\end{longtable}


{
\centering
\aicfig{media/inv_tikz.pdf}
}

\emph{Figure 8: Transistor schematic of the CMOS inverter, one PMOS to
the supply and one NMOS to ground sharing gate A and output Y}


{
\centering
\aicfig{media/pdpu_tikz.pdf}
}

\emph{Figure 9: Output state as a function of the pull-up and pull-down
networks - Z when both are off, 1 or 0 when one conducts, and X when
both conduct}

\emph{PD = Pull-down PU = Pull-up}

\begin{Shaded}
\begin{Highlighting}[]
\NormalTok{logic }\OperatorTok{=\textgreater{}} \OperatorTok{[}\DecValTok{0}\OperatorTok{,}\DecValTok{1}\OperatorTok{,}\NormalTok{Z}\OperatorTok{,}\NormalTok{X}\OperatorTok{];}
\end{Highlighting}
\end{Shaded}


{
\centering
\aicfig{media/pull_tikz.pdf}
}

\emph{Figure 10: Block view of a static CMOS gate, where the inputs
drive a PMOS pull-up network and an NMOS pull-down network that share
the output node}

\emph{Pull-up series}

\begin{longtable}[]{@{}lll@{}}
\toprule\noalign{}
A & B & Y \\
\midrule\noalign{}
\endhead
\bottomrule\noalign{}
\endlastfoot
0 & 0 & 1 \\
0 & 1 & Z \\
1 & 0 & Z \\
1 & 1 & Z \\
\end{longtable}

\emph{Pull-up parallel}

\begin{longtable}[]{@{}lll@{}}
\toprule\noalign{}
A & B & Y \\
\midrule\noalign{}
\endhead
\bottomrule\noalign{}
\endlastfoot
0 & 0 & 1 \\
0 & 1 & 1 \\
1 & 0 & 1 \\
1 & 1 & Z \\
\end{longtable}


{
\centering
\aicfig{media/pu_pmos_tikz.pdf}
}

\emph{Figure 11: PMOS pull-up networks - two devices in series conduct
only when both A and B are 0, two in parallel conduct unless both are 1}

\emph{Pull-down series}

\begin{longtable}[]{@{}lll@{}}
\toprule\noalign{}
A & B & Y \\
\midrule\noalign{}
\endhead
\bottomrule\noalign{}
\endlastfoot
0 & 0 & Z \\
0 & 1 & Z \\
1 & 0 & Z \\
1 & 1 & 0 \\
\end{longtable}

\emph{Pull-down parallel}

\begin{longtable}[]{@{}lll@{}}
\toprule\noalign{}
A & B & Y \\
\midrule\noalign{}
\endhead
\bottomrule\noalign{}
\endlastfoot
0 & 0 & Z \\
0 & 1 & 0 \\
1 & 0 & 0 \\
1 & 1 & 0 \\
\end{longtable}


{
\centering
\aicfig{media/pd_nmos_tikz.pdf}
}

\emph{Figure 12: NMOS pull-down networks - two devices in series conduct
only when both A and B are 1, two in parallel conduct when either is 1}

\subsection{Rules for inverting logic}\label{rules-for-inverting-logic}\index{rules for inverting logic@Rules for inverting logic}

\textbf{Pull-up} OR =\textgreater{} PMOS in series =\textgreater{} POS
AND =\textgreater{} PMOS in parallel =\textgreater{} PAP

\textbf{Pull-down} OR =\textgreater{} NMOS in parallel =\textgreater{}
NOP AND =\textgreater{} NMOS in series =\textgreater{} NAS


{
\centering
\aicfig{media/pull_tikz.pdf}
}

\emph{Figure 13: The same pull-up and pull-down blocks, read through the
rules that OR maps to PMOS in series and NMOS in parallel, AND to PMOS
in parallel and NMOS in series}

\[ \text{Y} = \overline{\text{AB}} = \text{NOT ( A AND B)}\]

\textbf{AND} PU =\textgreater{} PMOS in parallel PD =\textgreater{} NMOS
in series


{
\centering
\aicfig{media/nand_tr_tikz.pdf}
}

\emph{Figure 14: NAND built from the AND rule, with PMOS A and B in
parallel for the pull-up and NMOS A and B in series for the pull-down}

\begin{longtable}[]{@{}lll@{}}
\toprule\noalign{}
A & B & \emph{NOT(A AND B)} \\
\midrule\noalign{}
\endhead
\bottomrule\noalign{}
\endlastfoot
0 & 0 & 1 \\
0 & 1 & 1 \\
1 & 0 & 1 \\
1 & 1 & 0 \\
\end{longtable}


{
\centering
\aicfig{media/nand_tikz.pdf}
}

\emph{Figure 15: Schematic symbol for the two input NAND gate}

\[ \text{Y} = \overline{\text{A + B}} = \text{NOT ( A OR B)}\]

\textbf{OR} PU =\textgreater{} PMOS in series PD =\textgreater{} NMOS in
parallel


{
\centering
\aicfig{media/nor_tr_tikz.pdf}
}

\emph{Figure 16: NOR built from the OR rule, with PMOS A and B in series
for the pull-up and NMOS A and B in parallel for the pull-down}

\begin{longtable}[]{@{}lll@{}}
\toprule\noalign{}
A & B & \emph{NOT(A OR B)} \\
\midrule\noalign{}
\endhead
\bottomrule\noalign{}
\endlastfoot
0 & 0 & 1 \\
0 & 1 & 0 \\
1 & 0 & 0 \\
1 & 1 & 0 \\
\end{longtable}


{
\centering
\aicfig{media/nor_tikz.pdf}
}

\emph{Figure 17: Schematic symbol for the two input NOR gate}

\section{SR-Latch}\label{sr-latch}\index{sr-latch@SR-Latch}

Use boolean expressions to figure out how gates work.

Remember De-Morgan

\[\overline{AB}  = \overline{A}+ \overline{B}\]
\[\overline{A+B}  = \overline{A} \cdot \overline{B}\]

\[Q = \overline{R \overline{Q}} = \overline{R} +
\overline{\overline{Q}} = \overline{R} + Q \]

\[\overline{Q} = \overline{S Q} = \overline{S} +
\overline{Q} = \overline{S} + \overline{Q} \]


{
\centering
\aicfig{media/sr.pdf}
}

\emph{Figure 18: SR-latch truth table and symbol, and the cross-coupled
NAND implementation}

\[Q = \overline{R} + Q\] ,

\[\overline{Q} =\overline{S} + \overline{Q}\]

\begin{longtable}[]{@{}llll@{}}
\toprule\noalign{}
S & R & Q & \textasciitilde Q \\
\midrule\noalign{}
\endhead
\bottomrule\noalign{}
\endlastfoot
0 & 0 & X & X \\
0 & 1 & 0 & 1 \\
1 & 0 & 1 & 0 \\
1 & 1 & Q & \textasciitilde Q \\
\end{longtable}

\section{D-Latch (16 transistors)}\label{d-latch-16-transistors}\index{d-latch (16 transistors)@D-latch (16 transistors)}

\begin{longtable}[]{@{}llll@{}}
\toprule\noalign{}
C & D & Q & \textasciitilde Q \\
\midrule\noalign{}
\endhead
\bottomrule\noalign{}
\endlastfoot
0 & X & Q & \textasciitilde Q \\
1 & 0 & 0 & 1 \\
1 & 1 & 1 & 0 \\
\end{longtable}


{
\centering
\aicfig{media/dlatch_tikz.pdf}
}

\emph{Figure 19: Gate level D latch built from NAND gates, with the
cross-coupled NAND SR latch and its truth table above and the full latch
driven by clock C below}

\section{Other logic cells}\label{other-logic-cells}\index{other logic cells@Other logic cells}

What about \[\text{Y} = \text{AB}\] and
\[\text{Y} = \text{A} + \text{B}\]?

\[\text{Y} = \text{AB} = \overline{\overline{\text{AB}}}\]

\textbf{Y} = \textbf{A} AND \textbf{B} = NOT( NOT( \textbf{A} AND
\textbf{B} ) )


{
\centering
\aicfig{media/and_tikz.pdf}
}

\emph{Figure 20: An AND gate built as a NAND followed by an inverter}

\[\text{Y} = \text{A+B} = \overline{\overline{\text{A+B}}}\]

\textbf{Y} = \textbf{A} OR \textbf{B} = NOT( NOT( \textbf{A} OR
\textbf{B} ) )


{
\centering
\aicfig{media/or_tikz.pdf}
}

\emph{Figure 21: An OR gate built as a NOR followed by an inverter}

\section{AOI22: and or invert}\label{aoi22-and-or-invert}\index{aoi22: and or invert@AOI22: and or invert}

\textbf{Y} = NOT( \textbf{A} AND \textbf{B} OR \textbf{C} AND
\textbf{D})

\[\text{Y} =  \overline{\text{AB} + \text{CD}}\]


{
\centering
\aicfig{media/an2oi_tikz.pdf}
}

\emph{Figure 22: AOI22 gate at transistor level, with the
series-parallel pull-up network above the output and the pull-down
network below}


{
\centering
\aicfig{media/inv_tg_tikz.pdf}
}

\emph{Figure 23: An inverter followed by a transmission gate is
equivalent to a tristate inverter, redrawn as a single stack of four
series transistors}

\section{Tristate inverter}\label{tristate-inverter}\index{tristate inverter@Tristate inverter}

\begin{longtable}[]{@{}lll@{}}
\toprule\noalign{}
E & A & Y \\
\midrule\noalign{}
\endhead
\bottomrule\noalign{}
\endlastfoot
0 & 0 & Z \\
0 & 1 & Z \\
1 & 0 & 1 \\
1 & 1 & 0 \\
\end{longtable}


{
\centering
\aicfig{media/ivtrix_tikz.pdf}
}

\emph{Figure 24: Tristate inverter - symbol and the four transistor
stack where input A drives the outer devices and the enable E and its
complement drive the inner devices}

\section{Mux}\label{mux}\index{mux@Mux}

\begin{longtable}[]{@{}ll@{}}
\toprule\noalign{}
S & Y \\
\midrule\noalign{}
\endhead
\bottomrule\noalign{}
\endlastfoot
0 & NOT(P1) \\
1 & NOT(P0) \\
\end{longtable}


{
\centering
\aicfig{media/mux_tikz.pdf}
}

\emph{Figure 25: Two input multiplexer built from two tristate
inverters, with select S and its complement steering either P0 or P1 to
output Y}

D-Latch (12 transistors)


{
\centering
\aicfig{media/latch_tikz.pdf}
}

\emph{Figure 26: Twelve transistor D latch - input inverter, clocked
transmission gate and a feedback inverter pair holding the state, shown
as symbol, gate level and transistor level}

D-Flip Flop (\textless{} 26 transistors)


{
\centering
\aicfig{media/d_ff_tikz.pdf}
}

\emph{Figure 27: D flip-flop drawn as two D latches in series clocked on
opposite phases of C}


{
\centering
\aicfig{media/digital_ff_comb_tikz.pdf}
}

\emph{Figure 28: Synchronous digital design, with clocked flip-flop
banks at the input and output and a cloud of combinational logic between
them}

\section{There are other types of
logic}\label{there-are-other-types-of-logic}

\begin{itemize}
\tightlist
\item
  True single phase clock (TSPC) logic
\item
  Pass transistor logic
\item
  Transmission gate logic
\item
  Differential logic
\item
  Dynamic logic
\end{itemize}

Consider other types of logic ``rule breaking'', so you should know why
you need it.


{
\centering
\aicfig{media/fig_sar_logic.pdf}
}

\emph{Figure 29: Dynamic logic in a 9-bit SAR ADC - the binary weighted
capacitor array with per-bit logic slices above, and the transistor
level dynamic cells below}

\emph{\emph{Dynamic logic =\textgreater{} A Compiled 9-bit 20-MS/s
3.5-fJ/conv.step SAR ADC in 28-nm FDSOI for Bluetooth Low Energy
Receivers \citeproc{ref-wulff17}{{[}1{]}}}}

\section{Speed}\label{speed}\index{speed@Speed}


{
\centering
\aicfig{media/cpumax.pdf}
}

\emph{Figure 30: Maximum microprocessor clock frequency in GHz plotted
against year from 1971 to 2018, on a logarithmic axis}


{
\centering
\aicfig{media/digital_ff_comb_tikz.pdf}
}

\emph{Figure 31: The same flip-flop and combinational logic structure,
whose maximum clock rate is set by the delay through the logic between
two flip-flops}

\section{Flip-flops and speed}\label{flip-flops-and-speed}\index{flip-flops and speed@Flip-flops and speed}


{
\centering
\aicfig{media/d_ff_tikz.pdf}
}

\emph{Figure 32: D flip-flop as two cascaded latches, alongside the
SPICE subcircuit of the DFRNQNX1 standard cell that implements it}

\begin{Shaded}
\begin{Highlighting}[]
\NormalTok{dicex}\OperatorTok{/}\NormalTok{lib}\OperatorTok{/}\ConstantTok{SUN\_TR\_GF130N}\AttributeTok{.spi}\OperatorTok{:}

\AttributeTok{.}\ConstantTok{SUBCKT} \ConstantTok{DFRNQNX1\_CV} \ConstantTok{D} \ConstantTok{CK} \ConstantTok{RN} \ConstantTok{Q} \ConstantTok{QN} \ConstantTok{AVDD} \ConstantTok{AVSS}
\ConstantTok{XA0} \ConstantTok{AVDD} \ConstantTok{AVSS} \ConstantTok{TAPCELLB\_CV}
\ConstantTok{XA1} \ConstantTok{CK} \ConstantTok{RN} \ConstantTok{CKN} \ConstantTok{AVDD} \ConstantTok{AVSS} \ConstantTok{NDX1\_CV}
\ConstantTok{XA2} \ConstantTok{CKN} \ConstantTok{CKB} \ConstantTok{AVDD} \ConstantTok{AVSS} \ConstantTok{IVX1\_CV}
\ConstantTok{XA3} \ConstantTok{D} \ConstantTok{CKN} \ConstantTok{CKB} \ConstantTok{A0} \ConstantTok{AVDD} \ConstantTok{AVSS} \ConstantTok{IVTRIX1\_CV}
\ConstantTok{XA4} \ConstantTok{A1} \ConstantTok{CKB} \ConstantTok{CKN} \ConstantTok{A0} \ConstantTok{AVDD} \ConstantTok{AVSS} \ConstantTok{IVTRIX1\_CV}
\ConstantTok{XA5} \ConstantTok{A0} \ConstantTok{A1} \ConstantTok{AVDD} \ConstantTok{AVSS} \ConstantTok{IVX1\_CV}
\ConstantTok{XA6} \ConstantTok{A1} \ConstantTok{CKB} \ConstantTok{CKN} \ConstantTok{QN} \ConstantTok{AVDD} \ConstantTok{AVSS} \ConstantTok{IVTRIX1\_CV}
\ConstantTok{XA7} \ConstantTok{Q} \ConstantTok{CKN} \ConstantTok{CKB} \ConstantTok{RN} \ConstantTok{QN} \ConstantTok{AVDD} \ConstantTok{AVSS} \ConstantTok{NDTRIX1\_CV}
\ConstantTok{XA8} \ConstantTok{QN} \ConstantTok{Q} \ConstantTok{AVDD} \ConstantTok{AVSS} \ConstantTok{IVX1\_CV}
\AttributeTok{.}\ConstantTok{ENDS}
\end{Highlighting}
\end{Shaded}

Setup time: How long before clk does the data need to change

The setup time is not a number the flip-flop advertises, it is a number
you measure. Sweep the moment the data changes relative to the rising
clock edge, simulate, and look at where the output stops following the
input. The two plots below are two points either side of that boundary,
8 ps and 10 ps.


{
\centering
\aicfig{media/dff_setup_8_tikz.pdf}
}


{
\centering
\aicfig{media/dff_setup_10_tikz.pdf}
}

\emph{Figure 33: Simulated d, ck, q and qn of the D flip-flop for two
setup times, 8 ps and 10 ps. With 8 ps the data changes too close to the
rising clock edge at 0.5 ns, the flip-flop does not capture it, and q
only goes high at the second edge at 1.5 ns. With 10 ps the data has
settled early enough and q rises with the first edge. Two picoseconds
separate the two}

Left of that boundary the flip-flop still switches, but late: notice how
q in the 8 ps plot rises a full clock period after it should. That is
the failure mode setup violations produce in a real chip. Nothing is
stuck, nothing looks broken on a scope, the data simply arrives one
cycle behind, and it only happens on the corners and the temperatures
where the launching path is slowest.

Hold time: How long after clk can the data change

Hold time is the same experiment run from the other side. Now the
question is not whether the data arrived early enough to be captured,
but whether it stayed put long enough afterwards for the capture to
finish. Move the data edge towards the clock edge and the flip-flop
eventually samples the \emph{new} value instead of the old one.


{
\centering
\aicfig{media/dff_hold_-40_tikz.pdf}
}


{
\centering
\aicfig{media/dff_hold_-30_tikz.pdf}
}

\emph{Figure 34: The same signals for two hold times, -40 ps and -30 ps,
that is, the data changing 40 ps and 30 ps before the second rising
clock edge at 1.5 ns. At -40 ps the flip-flop takes the new low value
and q falls. At -30 ps the change is not taken and q stays high for
another period}

Hold violations are worse than setup violations, and it is worth being
clear about why. A setup violation you can fix after the fact by slowing
the clock down; the path simply needs more time, and the same silicon
works at a lower frequency. A hold violation does not care what the
clock frequency is. The data races the clock over a distance that has
nothing to do with the period, so a chip that fails hold fails at every
frequency, including DC, and the only fix is more silicon: buffers
inserted in the fast path, which means another place and route pass.

\section{Timing analysis}\label{timing-analysis}\index{timing analysis@Timing analysis}

Analyze arrival times of all nodes in a combinatorial circuit

\[ arrival_i = max_{j \in fanin(i)}{arrival_j} + t_{pd_i} \Rightarrow  a_i = max_{j \in fanin(i)}{a_j} + t_{pd_i}\]

\[ slack_i = required_i - arrival_i\]

Positive slack (over PVT\footnote{PVT =\textgreater{} Process, Voltage,
  Temperature}) =\textgreater{} Timing is OK Negative slack (over
PVT\footnote{PVT =\textgreater{} Process, Voltage, Temperature})
=\textgreater{} Timing is not OK


{
\centering
\aicfig{media/timing_tikz.pdf}
}

\emph{Figure 35: Arrival time propagation through a small combinational
circuit, where each gate adds its delay to the largest arrival time at
its inputs to give 130 at the output}

\section{Timing analysis tools}\label{timing-analysis-tools}\index{timing analysis tools@Timing analysis tools}

\textbf{Commercial}
\href{https://www.cadence.com/en_US/home/tools/digital-design-and-signoff/silicon-signoff/tempus-timing-signoff-solution.html}{Cadence
Tempus}

\href{https://www.synopsys.com/implementation-and-signoff/signoff/primetime.html}{Synopsys
PrimeTime}

\textbf{Free} \href{https://github.com/OpenTimer/OpenTimer}{OpenTimer}

\subsection{\texorpdfstring{\href{https://www.synopsys.com/glossary/what-is-static-timing-analysis.html}{What
is timing
analysis}}{What is timing analysis}}\label{what-is-timing-analysis}


{
\centering
\aicfig{media/timing_paths_tikz.pdf}
}

\emph{Figure 36: Timing paths through a circuit, each running from a
launch point through combinational logic to a capture point, with a
table of the four startpoint and endpoint combinations}


{
\centering
\aicfig{media/timing_types_tikz.pdf}
}

\emph{Figure 37: The path types considered in timing analysis - data
path, clock path, asynchronous reset path and clock-gating path}

\subsubsection{\texorpdfstring{\href{https://www.csee.umbc.edu/courses/graduate/CMPE641/Fall08/cpatel2/slides/lect05_LIB.pdf}{What
do the tools
need?}}{What do the tools need?}}\label{what-do-the-tools-need}

Input and output delay paths as a function of input transition time and
capacitive load, setup and hold time.

\href{https://github.com/OpenTimer/OpenTimer/blob/master/example/simple/osu018_stdcells.lib}{osu018\_stdcells.lib}

\begin{Shaded}
\begin{Highlighting}[]

\ErrorTok{cell} \ErrorTok{(INVX1)} \FunctionTok{\{}
  \ErrorTok{cell\_footprint} \FunctionTok{:} \ErrorTok{inv;}
\ErrorTok{area} \ErrorTok{:} \DecValTok{16}\ErrorTok{;}
  \ErrorTok{cell\_leakage\_power} \ErrorTok{:} \FloatTok{0.0221741}\ErrorTok{;}
  \ErrorTok{pin(A)}  \FunctionTok{\{}
    \ErrorTok{direction} \FunctionTok{:} \ErrorTok{input;}
    \ErrorTok{capacitance} \ErrorTok{:} \FloatTok{0.00932456}\ErrorTok{;}
    \ErrorTok{rise\_capacitance} \ErrorTok{:} \FloatTok{0.00932196}\ErrorTok{;}
    \ErrorTok{fall\_capacitance} \ErrorTok{:} \FloatTok{0.00932456}\ErrorTok{;}
  \FunctionTok{\}}
  \ErrorTok{pin(Y)}  \FunctionTok{\{}
    \ErrorTok{direction} \FunctionTok{:} \ErrorTok{output;}
    \ErrorTok{capacitance} \ErrorTok{:} \DecValTok{0}\ErrorTok{;}
    \ErrorTok{rise\_capacitance} \ErrorTok{:} \DecValTok{0}\ErrorTok{;}
    \ErrorTok{fall\_capacitance} \ErrorTok{:} \DecValTok{0}\ErrorTok{;}
    \ErrorTok{max\_capacitance} \ErrorTok{:} \FloatTok{0.503808}\ErrorTok{;}
    \ErrorTok{function} \ErrorTok{:} \StringTok{"(!A)"}\ErrorTok{;}
    \ErrorTok{timing()} \FunctionTok{\{}
      \ErrorTok{related\_pin} \FunctionTok{:} \StringTok{"A"}\ErrorTok{;}
      \ErrorTok{timing\_sense} \ErrorTok{:} \ErrorTok{negative\_unate;}
      \ErrorTok{cell\_fall(delay\_template\_}\DecValTok{5}\ErrorTok{x5)} \FunctionTok{\{}
        \ErrorTok{index\_1} \ErrorTok{(}\DataTypeTok{"0.005, 0.0125, 0.025, 0.075, 0.15"}\ErrorTok{);}
        \ErrorTok{index\_2} \ErrorTok{(}\DataTypeTok{"0.06, 0.18, 0.42, 0.6, 1.2"}\ErrorTok{);}
        \ErrorTok{values} \ErrorTok{(} \ErrorTok{\textbackslash{}}
          \DataTypeTok{"0.030906, 0.037434, 0.038584, 0.039088, 0.030318"}\FunctionTok{,} \ErrorTok{\textbackslash{}}
          \DataTypeTok{"0.04464, 0.057551, 0.073142, 0.077841, 0.081003"}\FunctionTok{,} \ErrorTok{\textbackslash{}}
          \DataTypeTok{"0.064368, 0.091076, 0.11557, 0.126352, 0.144944"}\FunctionTok{,} \ErrorTok{\textbackslash{}}
          \DataTypeTok{"0.139135, 0.174422, 0.232659, 0.261317, 0.321043"}\FunctionTok{,} \ErrorTok{\textbackslash{}}
          \DataTypeTok{"0.249412, 0.28434, 0.357694, 0.406534, 0.51187"}\ErrorTok{);}
      \FunctionTok{\}}
      
\end{Highlighting}
\end{Shaded}

\begin{Shaded}
\begin{Highlighting}[]
      \ErrorTok{fall\_transition(delay\_template\_5x5)} \FunctionTok{\{}
        \ErrorTok{index\_1} \ErrorTok{(}\DataTypeTok{"0.005, 0.0125, 0.025, 0.075, 0.15"}\ErrorTok{);}
        \ErrorTok{index\_2} \ErrorTok{(}\DataTypeTok{"0.06, 0.18, 0.42, 0.6, 1.2"}\ErrorTok{);}
        \ErrorTok{values} \ErrorTok{(} \ErrorTok{\textbackslash{}}
          \DataTypeTok{"0.032269, 0.0648, 0.087, 0.1032, 0.1476"}\FunctionTok{,} \ErrorTok{\textbackslash{}}
          \DataTypeTok{"0.036025, 0.0726, 0.1044, 0.1236, 0.183"}\FunctionTok{,} \ErrorTok{\textbackslash{}}
          \DataTypeTok{"0.06, 0.0882, 0.1314, 0.1554, 0.2286"}\FunctionTok{,} \ErrorTok{\textbackslash{}}
          \DataTypeTok{"0.1494, 0.1578, 0.2124, 0.2508, 0.3528"}\FunctionTok{,} \ErrorTok{\textbackslash{}}
          \DataTypeTok{"0.288, 0.2892, 0.3192, 0.3576, 0.492"}\ErrorTok{);}
      \FunctionTok{\}}
      \ErrorTok{cell\_rise(delay\_template\_5x5)} \FunctionTok{\{}
        \ErrorTok{index\_1} \ErrorTok{(}\DataTypeTok{"0.005, 0.0125, 0.025, 0.075, 0.15"}\ErrorTok{);}
        \ErrorTok{index\_2} \ErrorTok{(}\DataTypeTok{"0.06, 0.18, 0.42, 0.6, 1.2"}\ErrorTok{);}
        \ErrorTok{values} \ErrorTok{(} \ErrorTok{\textbackslash{}}
          \DataTypeTok{"0.037639, 0.056898, 0.083401, 0.104927, 0.156652"}\FunctionTok{,} \ErrorTok{\textbackslash{}}
          \DataTypeTok{"0.05258, 0.083003, 0.119028, 0.141927, 0.207952"}\FunctionTok{,} \ErrorTok{\textbackslash{}}
          \DataTypeTok{"0.07402, 0.112622, 0.162437, 0.191122, 0.271755"}\FunctionTok{,} \ErrorTok{\textbackslash{}}
          \DataTypeTok{"0.15767, 0.201007, 0.284096, 0.331746, 0.452958"}\FunctionTok{,} \ErrorTok{\textbackslash{}}
          \DataTypeTok{"0.285016, 0.326868, 0.415086, 0.481337, 0.653064"}\ErrorTok{);}
      \FunctionTok{\}}
      \ErrorTok{rise\_transition(delay\_template\_5x5)} \FunctionTok{\{}
        \ErrorTok{index\_1} \ErrorTok{(}\DataTypeTok{"0.005, 0.0125, 0.025, 0.075, 0.15"}\ErrorTok{);}
        \ErrorTok{index\_2} \ErrorTok{(}\DataTypeTok{"0.06, 0.18, 0.42, 0.6, 1.2"}\ErrorTok{);}
        \ErrorTok{values} \ErrorTok{(} \ErrorTok{\textbackslash{}}
          \DataTypeTok{"0.031447, 0.059488, 0.0846, 0.0918, 0.138"}\FunctionTok{,} \ErrorTok{\textbackslash{}}
          \DataTypeTok{"0.047167, 0.0786, 0.1044, 0.1224, 0.1734"}\FunctionTok{,} \ErrorTok{\textbackslash{}}
          \DataTypeTok{"0.072, 0.096, 0.1398, 0.1578, 0.222"}\FunctionTok{,} \ErrorTok{\textbackslash{}}
          \DataTypeTok{"0.1866, 0.1914, 0.2358, 0.2748, 0.3696"}\FunctionTok{,} \ErrorTok{\textbackslash{}}
          \DataTypeTok{"0.3648, 0.3648, 0.384, 0.4146, 0.5388"}\ErrorTok{);}
      \FunctionTok{\}}
    \ErrorTok{\}}
    \ErrorTok{internal\_power()} \FunctionTok{\{}
      \ErrorTok{related\_pin} \FunctionTok{:} \StringTok{"A"}\ErrorTok{;}
      \ErrorTok{fall\_power(energy\_template\_}\DecValTok{5}\ErrorTok{x5)} \FunctionTok{\{}
        \ErrorTok{index\_1} \ErrorTok{(}\DataTypeTok{"0.005, 0.0125, 0.025, 0.075, 0.15"}\ErrorTok{);}
        \ErrorTok{index\_2} \ErrorTok{(}\DataTypeTok{"0.06, 0.18, 0.42, 0.6, 1.2"}\ErrorTok{);}
        \ErrorTok{values} \ErrorTok{(} \ErrorTok{\textbackslash{}}
          \DataTypeTok{"0.009213, 0.004772, 0.00823, 0.018532, 0.054083"}\FunctionTok{,} \ErrorTok{\textbackslash{}}
          \DataTypeTok{"0.009047, 0.005677, 0.005713, 0.015244, 0.049453"}\FunctionTok{,} \ErrorTok{\textbackslash{}}
          \DataTypeTok{"0.008669, 0.006332, 0.002998, 0.01159, 0.04368"}\FunctionTok{,} \ErrorTok{\textbackslash{}}
          \DataTypeTok{"0.007879, 0.007243, 0.001451, 0.004701, 0.030385"}\FunctionTok{,} \ErrorTok{\textbackslash{}}
          \DataTypeTok{"0.007605, 0.007297, 0.003652, 0.000737, 0.020842"}\ErrorTok{);}
      \FunctionTok{\}}
      \ErrorTok{rise\_power(energy\_template\_}\DecValTok{5}\ErrorTok{x5)} \FunctionTok{\{}
        \ErrorTok{index\_1} \ErrorTok{(}\DataTypeTok{"0.005, 0.0125, 0.025, 0.075, 0.15"}\ErrorTok{);}
        \ErrorTok{index\_2} \ErrorTok{(}\DataTypeTok{"0.06, 0.18, 0.42, 0.6, 1.2"}\ErrorTok{);}
        \ErrorTok{values} \ErrorTok{(} \ErrorTok{\textbackslash{}}
          \DataTypeTok{"0.023555, 0.029044, 0.041387, 0.051684, 0.087278"}\FunctionTok{,} \ErrorTok{\textbackslash{}}
          \DataTypeTok{"0.023165, 0.028621, 0.039211, 0.048916, 0.083039"}\FunctionTok{,} \ErrorTok{\textbackslash{}}
          \DataTypeTok{"0.023574, 0.02752, 0.036904, 0.045723, 0.077971"}\FunctionTok{,} \ErrorTok{\textbackslash{}}
          \DataTypeTok{"0.024479, 0.025247, 0.032268, 0.039242, 0.066587"}\FunctionTok{,} \ErrorTok{\textbackslash{}}
          \DataTypeTok{"0.024942, 0.025187, 0.029612, 0.034835, 0.057524"}\ErrorTok{);}
      \FunctionTok{\}}
    \FunctionTok{\}}
  \ErrorTok{\}}
\ErrorTok{\}}
\end{Highlighting}
\end{Shaded}

\section{Every gate must be simulated to provide behavior over input
transition and load
capacitance}\label{every-gate-must-be-simulated-to-provide-behavior-over-input-transition-and-load-capacitance}

\section{\texorpdfstring{All analog blocks must have associated liberty
file to describe behavior and timing paths \emph{If you integrate analog
into digital top
flow}}{All analog blocks must have associated liberty file to describe behavior and timing paths If you integrate analog into digital top flow}}\label{all-analog-blocks-must-have-associated-liberty-file-to-describe-behavior-and-timing-paths-if-you-integrate-analog-into-digital-top-flow}

\section{Gate Delay}\label{gate-delay}\index{gate delay@Gate delay}

\subsection{Delay definitions}\label{delay-definitions}\index{delay definitions@Delay definitions}

\begin{longtable}[]{@{}
  >{\raggedright\arraybackslash}p{(\linewidth - 4\tabcolsep) * \real{0.3333}}
  >{\raggedright\arraybackslash}p{(\linewidth - 4\tabcolsep) * \real{0.3333}}
  >{\raggedright\arraybackslash}p{(\linewidth - 4\tabcolsep) * \real{0.3333}}@{}}
\toprule\noalign{}
\begin{minipage}[b]{\linewidth}\raggedright
Parameter
\end{minipage} & \begin{minipage}[b]{\linewidth}\raggedright
Name
\end{minipage} & \begin{minipage}[b]{\linewidth}\raggedright
Description
\end{minipage} \\
\midrule\noalign{}
\endhead
\bottomrule\noalign{}
\endlastfoot
t\_pdr & max rising propagation delay & input to rising output cross 50
\% \\
t\_pdf & max falling propagation delay & input to falling output cross
50 \% \\
t\_pd & propagation delay & t\_pd = (t\_pdr + t\_pdf)/2 \\
t\_r & rise time & 20 \% to 80 \% \\
\end{longtable}

\begin{longtable}[]{@{}
  >{\raggedright\arraybackslash}p{(\linewidth - 4\tabcolsep) * \real{0.3333}}
  >{\raggedright\arraybackslash}p{(\linewidth - 4\tabcolsep) * \real{0.3333}}
  >{\raggedright\arraybackslash}p{(\linewidth - 4\tabcolsep) * \real{0.3333}}@{}}
\toprule\noalign{}
\begin{minipage}[b]{\linewidth}\raggedright
Parameter
\end{minipage} & \begin{minipage}[b]{\linewidth}\raggedright
Name
\end{minipage} & \begin{minipage}[b]{\linewidth}\raggedright
Description
\end{minipage} \\
\midrule\noalign{}
\endhead
\bottomrule\noalign{}
\endlastfoot
t\_f & fall time & 80 \% to 20 \% \\
t\_cdr & min rising contamination delay & input to rising output cross
50 \% \\
t\_cdf & min falling contamination delay & input to falling output cross
50 \% \\
t\_cd & contamination delay & t\_cd = (t\_cdr + t\_cdf)/2 \\
\end{longtable}

\section{Delay estimation}\label{delay-estimation}\index{delay estimation@Delay estimation}

How can we get a reasonably accurate hand calculation model of delay?

\[ C \approx 1 \text{ fF}/\mu\text{m}\]

\[ R \approx 1 \text{ k}\Omega\mu\text{m}\]

\subsection{Inverter with inverter
load}\label{inverter-with-inverter-load}

\[ C \approx 1 \text{ fF}/\mu\text{m}\],
\[ R \approx 1 \text{ k}\Omega\mu\text{m}\]

\[ t_{pd} = R \times 6C  = 6RC \]

\[ t_{pd} = 6 \times 1 \times 10^{3} \times 1 \times 10^{-15} \text{ s}\]

\[ t_{pd} = 6 \times 10^{-12}  = 6 \text{ ps}\]

\section{Elmore Delay}\label{elmore-delay}\index{elmore delay@Elmore delay}

\[ t_{pd} \approx \sum_{\text{nodes}}{R_{\text{nodes}-to-source} C_i} \]

\[ = R_1C_1 + (R_1 + R_2)C_2 + ... + (R_1 + R_2 + ... + R_N) C_N\]

Good enough for hand calculation

\section{Delay components}\label{delay-components}\index{delay components@Delay components}

\textbf{Parasitic delay (p)}

p = 9 or 12 RC

Independent of load capacitance

\textbf{Effort delay (f)}

f = 5h RC

Proportional to load capacitance

Let's use process independent unit \[d = \frac{d_{real}}{\tau}\],
\[ \tau = 3 RC\]

Parasitic delay \[\Rightarrow p = 12 RC / 3 RC = 4\]

Effort delay \[\Rightarrow f = 5h RC / 3RC = \frac{5}{3} h \]

Delay \[\Rightarrow d = f + p = \frac{5}{3}h + 4\]

Logical effort (g) is the ratio of the input capacitance of a gate to
the input capacitance of an inverter delivering the same output current

Parasitic delay \[\Rightarrow p = 4\]

Logic effort \[\Rightarrow g = \frac{5}{3} \]

Electrical effort \[\Rightarrow h = 1\]

Effort \[\Rightarrow f = gh \]

Delay \[\Rightarrow d = f + p = gh + p = 5\frac{2}{3}\]

Real delay
\[\Rightarrow d = 5\frac{2}{3} \times 3 \text{ ps} = 17 \text{ ps}\]


{
\centering
\aicfig{media/fig_logeffort.pdf}
}

\emph{Figure 38: Logical effort summary table giving the stage and path
expressions for number of stages, logical effort, electrical effort,
branching effort, effort, effort delay, parasitic delay and total delay}

\section[Modern IC timing analysis requires computers with advanced
programs]{\texorpdfstring{Modern IC timing analysis requires computers
with advanced
programs\footnote{Opportunity for good programmers}}{Modern IC timing analysis requires computers with advanced programs}}\label{modern-ic-timing-analysis-requires-computers-with-advanced-programs2}

\section{Best number of stages}\label{best-number-of-stages}\index{best number of stages@Best number of stages}

\section{Which has shortest delay?}\label{which-has-shortest-delay}\index{which has shortest delay?@Which has shortest delay?}


{
\centering
\aicfig{media/path_tikz.pdf}
}

\emph{Figure 39: Two ways to drive a load of 64 - a single inverter, and
a chain of three inverters sized 1, 4 and 16, with the path effort delay
worked out for each}


{
\centering
\aicfig{media/fig_logeffort.pdf}
}

\emph{Figure 40: The logical effort table used to evaluate the two
candidate paths, giving H equal to 64, G equal to 1, B equal to 1 and
path effort F equal to 64}

\[H = C_{cout}/C_{in} = 64 \]

\[G = \prod{g_i} = \prod{1} = 1\]

\[B = 1\]

\[F = GBH = 64\]

\emph{One stage, with Sutherland's classic \[p \approx 1\] per inverter}
\[f = 64 \Rightarrow D = 64 + 1 = 65\]

\emph{Three stage with \[f=4\]}
\[D_F = 12, p = 3 \Rightarrow D = 12 + 3 = 15\]

The classic textbook numbers above use Sutherland's parasitic delay of
about 1 per inverter; our diffusion-heavy estimate earlier gave
\(p = 4\). Redo the sums with \(p = 4\) and you get 68 against 24 - the
absolute numbers move, the conclusion does not: split the path into
stages of effort around four.

\begin{center}\rule{0.5\linewidth}{0.5pt}\end{center}

For close to optimal delay, use \[f = 4\] \emph{}(Used to be
\[f=e\])\emph{}

\section{Trends}\label{trends}\index{trends@Trends}


{
\centering
\aicfig{media/rosc_vdd_tikz.pdf}
}

\emph{Figure 41: Ring oscillator frequency and power against supply
voltage, with the supply sensitivity and the energy per cycle below. The
sensitivity peaks near 0.6 V at more than 4 GHz per volt, and the energy
per cycle is worst where the ring is fastest}


{
\centering
\aicfig{media/rosc_temp_tikz.pdf}
}

\emph{Figure 42: Ring oscillator frequency falling from 2.14 GHz to 0.84
GHz, a factor of 2.6, as temperature rises from minus 40 to 150 degrees
Celsius. The slope below runs from about minus 12 MHz per degree in the
cold to minus 3 when hot, so the sensitivity is itself temperature
dependent}

\section{Attack vector}\label{attack-vector}\index{attack vector@Attack vector}

\begin{Shaded}
\begin{Highlighting}[]
\KeywordTok{module}\NormalTok{ counter}\OperatorTok{(}
               \DataTypeTok{output}\NormalTok{ logic }\OperatorTok{[}\NormalTok{WIDTH}\DecValTok{{-}1}\OperatorTok{:}\DecValTok{0}\OperatorTok{]}\NormalTok{ out}\OperatorTok{,}
               \DataTypeTok{input}\NormalTok{ logic              clk}\OperatorTok{,}
               \DataTypeTok{input}\NormalTok{ logic              reset}
               \OperatorTok{);}

   \DataTypeTok{parameter}\NormalTok{ WIDTH }\OperatorTok{=} \DecValTok{8}\OperatorTok{;}

\NormalTok{   logic }\OperatorTok{[}\NormalTok{WIDTH}\DecValTok{{-}1}\OperatorTok{:}\DecValTok{0}\OperatorTok{]}\NormalTok{                    count}\OperatorTok{;}
\NormalTok{   always\_comb }\KeywordTok{begin}
\NormalTok{      count }\OperatorTok{=}\NormalTok{ out }\OperatorTok{+} \DecValTok{1}\OperatorTok{;}
   \KeywordTok{end}

\NormalTok{   always\_ff }\OperatorTok{@(}\KeywordTok{posedge}\NormalTok{ clk }\DataTypeTok{or} \KeywordTok{posedge}\NormalTok{ reset}\OperatorTok{)} \KeywordTok{begin}
      \KeywordTok{if} \OperatorTok{(}\NormalTok{reset}\OperatorTok{)}
\NormalTok{        out }\OperatorTok{\textless{}=} \DecValTok{0}\OperatorTok{;}
      \KeywordTok{else}
\NormalTok{        out }\OperatorTok{\textless{}=}\NormalTok{ count}\OperatorTok{;}
   \KeywordTok{end}

\KeywordTok{endmodule} \CommentTok{// counter}
\end{Highlighting}
\end{Shaded}


{
\centering
\aicfig{media/counter_gtkw.png}
}

\emph{Figure 43: GTKWave view of the 8-bit counter testbench, where the
count bus ramps over 2.6 us while clk toggles and reset stays low}


{
\centering
\aicfig{media/counter_gf130n.png}
}

\emph{Figure 44: The same 8-bit counter synthesized to the SUN TR GF130N
standard cell library, with eight D flip-flops on the right and the
increment logic to the left}

\begin{verbatim}
.SUBCKT counter out_7 out_6 out_5 out_4 out_3 out_2 out_1 out_0 clk reset AVDD AVSS
* SPICE netlist generated by Yosys 0_9 (git sha1 1979e0b1, gcc 10_3_0-1ubuntu1~20_10 -fPIC -Os)
X0 out_2 1 AVDD AVSS IVX1_CV
X1 out_3 2 AVDD AVSS IVX1_CV
X2 out_4 3 AVDD AVSS IVX1_CV
X3 out_5 4 AVDD AVSS IVX1_CV
X4 out_6 5 AVDD AVSS IVX1_CV
X5 out_0 6 AVDD AVSS IVX1_CV
X6 out_1 7 AVDD AVSS IVX1_CV
X7 6 7 8 AVDD AVSS NRX1_CV
X8 out_0 out_1 9 AVDD AVSS NDX1_CV
X9 1 9 10 AVDD AVSS NRX1_CV
X10 10 11 AVDD AVSS IVX1_CV
X11 2 11 12 AVDD AVSS NRX1_CV
X12 out_3 10 13 AVDD AVSS NDX1_CV
X13 out_3 10 14 AVDD AVSS NRX1_CV
X14 12 14 15 AVDD AVSS NRX1_CV
X15 3 13 16 AVDD AVSS NRX1_CV
X16 16 17 AVDD AVSS IVX1_CV
X17 out_4 12 18 AVDD AVSS NRX1_CV
X18 16 18 19 AVDD AVSS NRX1_CV
X19 4 17 20 AVDD AVSS NRX1_CV
X20 out_5 16 21 AVDD AVSS NDX1_CV
X21 out_5 16 22 AVDD AVSS NRX1_CV
X22 20 22 23 AVDD AVSS NRX1_CV
X23 5 21 24 AVDD AVSS NRX1_CV
X24 out_6 20 25 AVDD AVSS NRX1_CV
X25 24 25 26 AVDD AVSS NRX1_CV
X26 out_7 24 27 AVDD AVSS NRX1_CV
X27 out_7 24 28 AVDD AVSS NDX1_CV
X28 28 29 AVDD AVSS IVX1_CV
X29 27 29 30 AVDD AVSS NRX1_CV
X30 out_0 out_1 31 AVDD AVSS NRX1_CV
X31 8 31 32 AVDD AVSS NRX1_CV
X32 out_2 8 33 AVDD AVSS NRX1_CV
X33 10 33 34 AVDD AVSS NRX1_CV
X34 35 clk AVSS reset out_0 35 AVDD AVSS DFSRQNX1_CV
X35 32 clk AVSS reset out_1 36 AVDD AVSS DFSRQNX1_CV
X36 34 clk AVSS reset out_2 37 AVDD AVSS DFSRQNX1_CV
X37 15 clk AVSS reset out_3 38 AVDD AVSS DFSRQNX1_CV
X38 19 clk AVSS reset out_4 39 AVDD AVSS DFSRQNX1_CV
X39 23 clk AVSS reset out_5 40 AVDD AVSS DFSRQNX1_CV
X40 26 clk AVSS reset out_6 41 AVDD AVSS DFSRQNX1_CV
X41 30 clk AVSS reset out_7 42 AVDD AVSS DFSRQNX1_CV
V0 count_0 35 DC 0
V1 43 out_2 DC 0
V2 44 out_3 DC 0
V3 count_3 15 DC 0
V4 45 out_4 DC 0
V5 count_4 19 DC 0
V6 46 out_5 DC 0
V7 count_5 23 DC 0
V8 47 out_6 DC 0
V9 count_6 26 DC 0
V10 48 out_7 DC 0
V11 count_7 30 DC 0
V12 49 out_0 DC 0
V13 50 out_1 DC 0
V14 count_1 32 DC 0
V15 count_2 34 DC 0
.ENDS
\end{verbatim}


{
\centering
\aicfig{media/counter_ref_do.pdf}
}

\emph{Figure 45: ngspice transient of the reference counter output dor,
counting linearly from 0 to 255 over 128 ns}

dicex/sim/verilog/counter\_sv/counter\_attack\_tb.cir

\begin{Shaded}
\begin{Highlighting}[]
\ConstantTok{VDDA} \ConstantTok{AVDD\_ATTACK} \DecValTok{0}\NormalTok{ dc }\FloatTok{0.5}\NormalTok{ pulse(}\FloatTok{1.5} \FloatTok{0.6}\NormalTok{ tcd trf trf tapw taper)}
\end{Highlighting}
\end{Shaded}


{
\centering
\aicfig{media/counter_ref_wave.pdf}
}

\emph{Figure 46: ngspice transient of the attacked supply, pulsed
between 1.5 V and 0.6 V, together with the least significant counter
bit}


{
\centering
\aicfig{media/counter_ref_dor.pdf}
}

\emph{Figure 47: The attacked counter output do in red against the
reference output dor in blue, showing the counts lost each time the
supply is glitched}


{
\centering
\aicfig{media/chipwisperer.png}
}

\emph{Figure 48: The ChipWhisperer web page, an open source platform for
side channel power analysis and fault injection attacks on embedded
systems}

\section{Pick two}\label{pick-two}\index{pick two@Pick two}


{
\centering
\aicfig{media/optimization_tikz.pdf}
}

\emph{Figure 49: The design trade-off triangle between power, speed and
area or cost}

\section{Power}\label{power}\index{power@Power}

\section{What is power?}\label{what-is-power}\index{what is power?@What is power?}

Instantanious power: \[ P(t) = I(t)V(t)\]

Energy : \[ \int_0^T{P(t)dt} \] {[}J{]}

Average power: \[\frac{1}{T} \int_0^T{P(t)dt} \] {[}W or J/s{]}

\section{Power dissipated in a
resistor}\label{power-dissipated-in-a-resistor}

Ohm's Law \[V_R = I_R R\]

\[P_R = V_R I_R =  I_R^2 R  = \frac{V_R^2}{R} \]

\section{Charging a capacitor to VDD}\label{charging-a-capacitor-to-vdd}\index{charging a capacitor to vdd@Charging a capacitor to VDD}

Capacitor differential equation \[ I_C = C\frac{dV}{dt}\]

\[E_{C}  = \int_0^\infty{I_C V_C dt} = \int_0^\infty{ C \frac{dV}{dt} V_C dt} = \int_0^{V_C}{C V dV} = C\left[\frac{V^2}{2}\right]_0^{V_{DD}} \]

\[E_{C} = \frac{1}{2} C V_{DD}^2\]

\section{Energy to charge a capacitor to a voltage
VDD}\label{energy-to-charge-a-capacitor-to-a-voltage-vdd}

\[E_{C} = \frac{1}{2} C V_{DD}^2\]

\[I_{VDD} = I_C = C \frac{dV}{dt}\]

\[E_{VDD} = \int_0^\infty{I_{VDD} V_{DD} dt} = \int_0^\infty{C \frac{dV}{dt} V_{DD} dt} = C V_{DD}\int_0^{V_{DD}}{dV} = C V_{DD}^2\]

Only half the energy is stored on the capacitor, the rest is dissipated
in the PMOS

\section{Discharging a capacitor to
0}\label{discharging-a-capacitor-to-0}

\[E_{C} = \frac{1}{2} C V_{DD}^2\]

Voltage is pulled to ground, and the power is dissipated in the NMOS

\section{Power consumption of digital
circuits}\label{power-consumption-of-digital-circuits}

\[E_{VDD} = C V_{DD}^2\]

In a clock distribution network (chain of inverters), every output is
charged once per clock cycle

\[P_{VDD} = C V_{DD}^2 f\]

\section{Sources of power dissipation in CMOS
logic}\label{sources-of-power-dissipation-in-cmos-logic}

\[P_{total} = P_{dynamic} + P_{static}\]

\textbf{Dynamic power dissipation}

Charging and discharging load capacitances

\emph{short-circut} current, when PMOS and NMOS conduct at the same time

\[P_{dynamic} = P_{switching} + P_{short circuit}\]

\textbf{Static power dissipation}

Subthreshold leakage in OFF transistors

Gate leakage (tunneling current) through gate dielectric

Source/drain reverse bias PN junction leakage

\[P_{static} = \left( I_{sub} + I_{gate} + I_{pn} \right) V_{DD}\]

\section{Switching Power in logic
gates}\label{switching-power-in-logic-gates}

Only output node transitions from low to high consume power from
\[V_{DD}\]

Define \[P_i\] to be the probability that a node is 1

Define \[\overline{P_i} = 1 - P_i\] to be the probability that a node is
0

Define \textbf{activity factor (\[\alpha_i\])} as the
\textbf{probability of switching a node from 0 to 1}

If the probabilty is uncorrelated from cycle to cycle

\[\alpha_i = \overline{P_i}P_i\]

\section{Switching probability}\label{switching-probability}\index{switching probability@Switching probability}

Random data \[P = 0.5\], \[\alpha = 0.25\]

Clocks \[\alpha = 1\]


{
\centering
\aicfig{media/tb_sw_prob.pdf}
}

\emph{Figure 50: Probability that the output of an AND2, OR2, NAND2,
NOR2 or XOR2 gate is 1, expressed in the input probabilities PA and PB}


{
\centering
\aicfig{media/tb_sw_prob.pdf}
}

\emph{Figure 51: The gate output probability table applied to the
network below, where each NAND2 output has probability 3 over 4 of being
1}


{
\centering
\aicfig{media/prob_tikz.pdf}
}

\emph{Figure 52: Two NAND2 gates driving a NOR2 gate, giving output
probability PY of 1 over 16 and activity factor 15 over 256 when all
inputs have probability one half}

Assume \[P = P_A = P_B = P_C = P_D = \frac{1}{2}\]

\[P_X = P_Z =  1 - P P = 1 - \frac{1}{4} = \frac{3}{4}\]

\[\overline{P_X} = \overline{P_Y} = \frac{1}{4}\]

\[P_Y = \frac{1}{4} \times \frac{1}{4} = \frac{1}{16}\]

\[\alpha = \frac{1}{16}\left(1 - \frac{1}{16}\right) = \frac{15}{16}\frac{1}{16} = \frac{15}{256}\]


{
\centering
\aicfig{media/tb_sw_prob.pdf}
}

\emph{Figure 53: The same gate output probability table, used here
alongside the De Morgan simplification of the network}


{
\centering
\aicfig{media/prob_tikz.pdf}
}

\emph{Figure 54: The same NAND-NAND-NOR network, where De Morgan reduces
the function to ABCD so that PY is the product of the four input
probabilities, 1 over 16}

\[ \overline{\overline{\text{AB}} + \overline{\text{CD}}} \]

Use \emph{De Morgan} first
\[\overline{A+B}  = \overline{A} \cdot \overline{B}\]

\[\overline{\overline{\text{AB}} + \overline{\text{CD}}} = \overline{\overline{\text{AB}}} \overline{\overline{\text{CD}}} = ABCD\]

\[\Rightarrow P_Y = P_A P_B P_C P_D = \left(\frac{1}{2}\right)^4 = \frac{1}{16} \]

\[P_{tot} = \alpha C V_{DD}^2 f\]

\section{Strategies to reduce dynamic
power}\label{strategies-to-reduce-dynamic-power}

\begin{enumerate}
\def\labelenumi{\arabic{enumi}.}
\tightlist
\item
  Stop clock
\item
  Stop activity
\item
  Reduce clock frequency
\item
  Turn off VDD
\item
  Reduce VDD
\end{enumerate}


{
\centering
\aicfig{media/digital_ff_comb_tikz.pdf}
}

\emph{Figure 55: Synchronous logic drawn as banks of D flip-flops
separated by a combinational cloud, the switching activity that costs
the dynamic power}

\subsection[Stop clock ]{\texorpdfstring{Stop clock
\footnote{Often called \emph{clock gating}}}{Stop clock }}\label{stop-clock-3}


{
\centering
\aicfig{media/stop_clock_tikz.pdf}
}

\emph{Figure 56: Clock gating, where enable logic feeds a latch and an
AND gate so that the output clock only toggles the flip-flops and
combinational cloud when the block is enabled}

\subsection{Stop activity}\label{stop-activity}\index{stop activity@Stop activity}

Stopping the clock stops the flip-flops, but it does not necessarily
stop the combinational cloud. If the inputs to the cloud keep moving,
every gate inside it keeps charging and discharging its load, and the
\(\alpha C V_{DD}^2 f\) bill arrives whether or not anything downstream
cares about the answer. Stopping the activity means breaking the data
path into the cloud, so its inputs are held constant and the logic
inside simply stops toggling.


{
\centering
\aicfig{media/logic.pdf}
}


{
\centering
\aicfig{media/stop_activity.pdf}
}

\emph{Figure 57: The same flip-flop banks and combinational cloud drawn
twice. In the first drawing the cloud is marked at the points where the
dynamic power equation can be attacked. In the second the data path from
the first flip-flop bank into the cloud is broken, so the cloud sees a
constant input and stops switching, while the flip-flops still receive
their clock}

\subsection{Reduce frequency}\label{reduce-frequency}\index{reduce frequency@Reduce frequency}


{
\centering
\aicfig{media/reduce_freq.pdf}
}

\emph{Figure 58: Reducing clock frequency by clocking the big
combinational cloud with ClkB and the small cloud with the faster ClkA}

\subsection[Turn off power supply ]{\texorpdfstring{Turn off power
supply
\footnote{Often called power gating}}{Turn off power supply }}\label{turn-off-power-supply-4}


{
\centering
\aicfig{media/powergate.pdf}
}

\emph{Figure 59: Power gating with PWRUP header switches above the gated
logic block and AND gates holding the outputs}

\subsubsection{Reduce power supply}\label{reduce-power-supply}


{
\centering
\aicfig{media/reduce_vdd.pdf}
}

\emph{Figure 60: Two supply domains, fast logic on VDDH and slow logic
on VDDL, joined by a cross-coupled level shifter}

\subsubsection{Energy-Delay Product}\label{energy-delay-product}

\[ EDP = k\frac{C^2 V_{DD}^3}{(V_{DD}- V_t)^{\text{1 to 2}}}\]

Differentiating with respect to \[V_{DD}\] and setting the result to
\[0\] it's possible to work out that

\[ V_{DD-opt} = \frac{3}{3-\text{1 to 2}}V_t  \in[1.5,3]V_{t}\]

\section{Wires}\label{wires}\index{wires@Wires}

\section{Wire geometry}\label{wire-geometry}\index{wire geometry@Wire geometry}

Pitch = w + s

Aspect ratio (AR) = t/w

These days \[AR \approx 2\]

\section{Metal stack}\label{metal-stack}\index{metal stack@Metal stack}

Often 5 - 10 layers of metal

\begin{longtable}[]{@{}
  >{\centering\arraybackslash}p{(\linewidth - 6\tabcolsep) * \real{0.2500}}
  >{\centering\arraybackslash}p{(\linewidth - 6\tabcolsep) * \real{0.2500}}
  >{\centering\arraybackslash}p{(\linewidth - 6\tabcolsep) * \real{0.2500}}
  >{\centering\arraybackslash}p{(\linewidth - 6\tabcolsep) * \real{0.2500}}@{}}
\toprule\noalign{}
\begin{minipage}[b]{\linewidth}\centering
Metal
\end{minipage} & \begin{minipage}[b]{\linewidth}\centering
Material
\end{minipage} & \begin{minipage}[b]{\linewidth}\centering
Thickness
\end{minipage} & \begin{minipage}[b]{\linewidth}\centering
Purpose
\end{minipage} \\
\midrule\noalign{}
\endhead
\bottomrule\noalign{}
\endlastfoot
Metal 1 & Copper & Thin & in gate routing \\
Metal 3 - 5 & Copper & Thicker & Between gates routing \\
RDL & Aluminium & Ultra thick & Can tolerate high forces during wire
bonding. \\
\end{longtable}


{
\centering
\aicfig{media/skymetal.png}
}

\emph{Figure 61: Cross section of the Skywater metal stack, five copper
layers M1 to M5 below two aluminium layers M6 and M7}

\section{Metal routing rules on IC}\label{metal-routing-rules-on-ic}\index{metal routing rules on ic@Metal routing rules on IC}

Odd numbers metals =\textgreater{} Horizontal routing (as far as
possible)

Even numbers metals =\textgreater{} Vertical routing (as far as
possible)

\section{Modeling Interconnect}\label{modeling-interconnect}\index{modeling interconnect@Modeling interconnect}

\textbf{Resistance} narrow size impedes flow

\textbf{Capacitance} through under the leaky pipes

\textbf{Inductance} paddle wheel intertia opposes changes in flow rate

\section{Lumped model}\label{lumped-model}\index{lumped model@Lumped model}

Use 1-segment \[\pi\]-model for Elmore delay

\begin{Shaded}
\begin{Highlighting}[]
\ExtensionTok{C/2}\NormalTok{   R   C/2}
\ExtensionTok{{-}{-}{-}/\textbackslash{}/\textbackslash{}/\textbackslash{}{-}{-}{-}}
 \KeywordTok{|}        \KeywordTok{|}
\ExtensionTok{{-}{-}{-}}      \AttributeTok{{-}{-}{-}}
\ExtensionTok{{-}{-}{-}}      \AttributeTok{{-}{-}{-}} 
 \KeywordTok{|}        \KeywordTok{|}
\ExtensionTok{{-}{-}{-}}      \AttributeTok{{-}{-}{-}} 
 \ExtensionTok{{-}}        \AttributeTok{{-}}
\end{Highlighting}
\end{Shaded}

\section{Wire resistance}\label{wire-resistance}\index{wire resistance@Wire resistance}

\[ \text{resistivity} \Rightarrow \rho \text{ [} \Omega\text{m]} \]

\[ R = \frac{\rho}{t}\frac{l}{w} = R_\square \frac{l}{w} \]

\[ R_\square  = \text{sheet resistance [} \Omega/\square \text{]} \]

To find resistance, count the number of squares

\[ R = R_\square \times \text{\# of squares} \]

\section{Most wires: Copper}\label{most-wires-copper}\index{most wires: copper@Most wires: copper}

\[R_{sheet-m1} \approx \frac{1.7 \mu\Omega cm}{200 nm} \approx 0.1 \Omega/\square\]\\
\[R_{sheet-m9} \approx \frac{1.7 \mu\Omega cm}{3 \mu m} \approx 0.006 \Omega/\square\]

\textbf{Pitfalls}

Cu atoms diffuse into silicon and can cause damage

Must be surrounded by a diffusion barrier

Difficult high current densities (mA/\[\mu\]m) and high temperature (125
C)


{
\centering
\aicfig{media/metals.png}
}

\emph{Figure 62: Bulk resistivity in micro-ohm cm for silver, copper,
gold, aluminium, tungsten and titanium}

\section{Contacts}\label{contacts}

Contacts and vias can have 2-20 \[\Omega\]

Must use many contacts/vias for high current wires

\section{Wire Capacitance}\label{wire-capacitance}\index{wire capacitance@Wire capacitance}

Dense wires has about \[0.2 \text{ fF/}\mu\text{m}\]

\section{FSM}\label{fsm}\index{fsm@FSM}

\section{Mealy machine}\label{mealy-machine}\index{mealy machine@Mealy machine}

An FSM where outputs depend on current state and inputs


{
\centering
\aicfig{media/mealy_machine_tikz.pdf}
}

\emph{Figure 63: Mealy machine, where the input drives both the
next-state logic and the output combinational block, so the output
depends on input and current state}

\section{Moore machine}\label{moore-machine}\index{moore machine@Moore machine}

An FSM where outputs depend on current state


{
\centering
\aicfig{media/moore_machine_tikz.pdf}
}

\emph{Figure 64: Moore machine, where the output combinational block is
driven only by the state register}

\section{Mealy versus Moore}\label{mealy-versus-moore}\index{mealy versus moore@Mealy versus Moore}

\begin{longtable}[]{@{}
  >{\centering\arraybackslash}p{(\linewidth - 4\tabcolsep) * \real{0.3333}}
  >{\centering\arraybackslash}p{(\linewidth - 4\tabcolsep) * \real{0.3333}}
  >{\centering\arraybackslash}p{(\linewidth - 4\tabcolsep) * \real{0.3333}}@{}}
\toprule\noalign{}
\begin{minipage}[b]{\linewidth}\centering
Parameter
\end{minipage} & \begin{minipage}[b]{\linewidth}\centering
Mealy
\end{minipage} & \begin{minipage}[b]{\linewidth}\centering
Moore
\end{minipage} \\
\midrule\noalign{}
\endhead
\bottomrule\noalign{}
\endlastfoot
Outputs & depend on input and current state & output depend on current
state \\
States & Same, or fewer states than Moore & \\
Inputs & React faster to inputs & Next clock cycle \\
Outputs & Can be asynchronous & Synchronous \\
States & Generally requires fewer states for synthesis & More states
than Mealy \\
Counter & A counter is not a mealy machine & A counter is a Moore
machine \\
Design & Can be tricky to design & Easy \\
\end{longtable}

\subsection{dicex/sim/counter\_sv/counter.v}\label{dicexsimcounter_svcounter.v}

\begin{Shaded}
\begin{Highlighting}[]
\KeywordTok{module}\NormalTok{ counter}\OperatorTok{(}
               \DataTypeTok{output}\NormalTok{ logic }\OperatorTok{[}\NormalTok{WIDTH}\DecValTok{{-}1}\OperatorTok{:}\DecValTok{0}\OperatorTok{]}\NormalTok{ out}\OperatorTok{,}
               \DataTypeTok{input}\NormalTok{ logic              clk}\OperatorTok{,}
               \DataTypeTok{input}\NormalTok{ logic              reset}
               \OperatorTok{);}
   \DataTypeTok{parameter}\NormalTok{ WIDTH                      }\OperatorTok{=} \DecValTok{8}\OperatorTok{;}
\NormalTok{   logic }\OperatorTok{[}\NormalTok{WIDTH}\DecValTok{{-}1}\OperatorTok{:}\DecValTok{0}\OperatorTok{]}\NormalTok{                    count}\OperatorTok{;}
   
\NormalTok{   always\_comb }\KeywordTok{begin}
\NormalTok{      count }\OperatorTok{=}\NormalTok{ out }\OperatorTok{+} \DecValTok{1}\OperatorTok{;}
   \KeywordTok{end}

\NormalTok{   always\_ff }\OperatorTok{@(}\KeywordTok{posedge}\NormalTok{ clk }\DataTypeTok{or} \KeywordTok{posedge}\NormalTok{ reset}\OperatorTok{)} \KeywordTok{begin}
      \KeywordTok{if} \OperatorTok{(}\NormalTok{reset}\OperatorTok{)}
\NormalTok{        out }\OperatorTok{\textless{}=} \DecValTok{0}\OperatorTok{;}
      \KeywordTok{else}
\NormalTok{        out }\OperatorTok{\textless{}=}\NormalTok{ count}\OperatorTok{;}
   \KeywordTok{end}

\KeywordTok{endmodule} \CommentTok{// counter}
\end{Highlighting}
\end{Shaded}

\section{Battery charger FSM}\label{battery-charger-fsm}\index{battery charger fsm@Battery charger FSM}


{
\centering
\aicfig{media/charge_graph_tikz.pdf}
}

\emph{Figure 65: Li-Ion charging profile, charge current and battery
voltage against time through trickle charge, fast charge, constant
voltage charge and charging complete}

\subsection{Li-Ion batteries}\label{li-ion-batteries}\index{li-ion batteries@Li-Ion batteries}

Most Li-Ion batteries can tolerate 1 C during fast charge

For Biltema 18650 cells: \[ 1\text{ C} = 2950\text{ mA}\]
\[ 0.1\text{ C} = 295\text{ mA}\]

Most Li-Ion need to be charged to a termination voltage of 4.2 V


{
\centering
\aicfig{media/18650.jpeg}
}

\emph{Figure 66: A Biltema ICR18650 rechargeable Li-ion cell rated 2950
mAh at 3.7 V}

\textbf{Too high termination voltage, or too high charging current can
cause growth of lithium dendrites, that short + and -. Will end in
flames. Always check manufacturer datasheet for charging curves and
voltages}

\subsection{Battery charger - Inputs}\label{battery-charger---inputs}\index{battery charger - inputs@Battery charger - Inputs}

Voltage above \[V_{TRICKLE}\]

Voltage close to \[V_{TERM}\]

If voltage close to \[V_{TERM}\] and current is close to \[I_{TERM}\],
then charging complete

If charging complete, and voltage has dropped (\[V_{RECHARGE}\]), then
start again


{
\centering
\aicfig{media/charge_graph_tikz.pdf}
}

\emph{Figure 67: The charging profile marked with the thresholds the
charger senses - the trickle to fast voltage, the termination voltage
VTERM and the termination current ITERM}

\subsection{Battery charger - States}\label{battery-charger---states}\index{battery charger - states@Battery charger - States}

Trickle charge (0.1 C)

Fast charge (1 C)

Constant voltage

Charging complete


{
\centering
\aicfig{media/charge_graph_tikz.pdf}
}

\emph{Figure 68: The charging profile divided into the four charger
states - trickle charge at 0.1 C, fast charge at 1 C, constant voltage
and charging complete}


{
\centering
\aicfig{media/bcharger.pdf}
}

\emph{Figure 69: Battery charger state machine, cycling trickle charge,
fast charge, constant voltage and complete on the vtrkl, vterm, iterm
and vrchrg flags}

\subsubsection{One way to draw FSMs -
Graphviz}\label{one-way-to-draw-fsms---graphviz}

\begin{verbatim}
digraph finite_state_machine {
    rankdir=LR;
    size="8,5"

    node [shape = doublecircle, label="Trickle charger", fontsize=12] trkl;
    node [shape = circle, label="Fast charge", fontsize=12] fast;
    node [shape = circle, label="Const. Voltage", fontsize=12] vconst;
    node [shape = circle, label="Done", fontsize=12] done;

    trkl -> trkl [label="vtrkl = 0"];
    trkl -> fast [label="vtrkl = 1"];
    fast -> fast [label="vterm = 0"];
    fast -> vconst [label="vterm = 1"];
    vconst-> vconst [label="iterm = 0"];
    vconst-> done [label="iterm = 1"];
    done-> done [label="vrchrg = 0"];
    done-> trkl [label="vrchrg = 1"];

}
\end{verbatim}

\begin{verbatim}
dot -Tpdf bcharger.dot -o bcharger.pdf
\end{verbatim}


{
\centering
\aicfig{media/bcharger.pdf}
}

\emph{Figure 70: The battery charger state diagram beside the
SystemVerilog case statement that implements the next-state logic}

\begin{Shaded}
\begin{Highlighting}[]
\KeywordTok{module}\NormalTok{ bcharger}\OperatorTok{(} \DataTypeTok{output}\NormalTok{ logic trkl}\OperatorTok{,}
        \DataTypeTok{output}\NormalTok{ logic fast}\OperatorTok{,} 
        \DataTypeTok{output}\NormalTok{ logic vconst}\OperatorTok{,}
        \DataTypeTok{output}\NormalTok{ logic done}\OperatorTok{,}
        \DataTypeTok{input}\NormalTok{ logic  vtrkl}\OperatorTok{,} 
        \DataTypeTok{input}\NormalTok{ logic  vterm}\OperatorTok{,} 
        \DataTypeTok{input}\NormalTok{ logic  iterm}\OperatorTok{,} 
        \DataTypeTok{input}\NormalTok{ logic  vrchrg}\OperatorTok{,}
        \DataTypeTok{input}\NormalTok{ logic  clk}\OperatorTok{,} 
        \DataTypeTok{input}\NormalTok{ logic  reset}
                    \OperatorTok{);}

   \DataTypeTok{parameter}\NormalTok{ TRLK }\OperatorTok{=} \DecValTok{0}\OperatorTok{,}\NormalTok{ FAST }\OperatorTok{=} \DecValTok{1}\OperatorTok{,}\NormalTok{ VCONST }\OperatorTok{=} \DecValTok{2}\OperatorTok{,}\NormalTok{ DONE}\OperatorTok{=}\DecValTok{3}\OperatorTok{;}
\NormalTok{   logic }\OperatorTok{[}\DecValTok{1}\OperatorTok{:}\DecValTok{0}\OperatorTok{]}\NormalTok{                   state}\OperatorTok{;}
\NormalTok{   logic }\OperatorTok{[}\DecValTok{1}\OperatorTok{:}\DecValTok{0}\OperatorTok{]}\NormalTok{                   next\_state}\OperatorTok{;}

   \CommentTok{//{-} Figure out the next state}
\NormalTok{   always\_comb }\KeywordTok{begin}
      \KeywordTok{case} \OperatorTok{(}\NormalTok{state}\OperatorTok{)}
        \DecValTok{TRLK:}\NormalTok{ next\_state }\OperatorTok{=}\NormalTok{ vtrkl }\OperatorTok{?}\NormalTok{ FAST }\OperatorTok{:}\NormalTok{ TRLK}\OperatorTok{;}
        \DecValTok{FAST:}\NormalTok{ next\_state }\OperatorTok{=}\NormalTok{ vterm }\OperatorTok{?}\NormalTok{ VCONST }\OperatorTok{:}\NormalTok{ FAST}\OperatorTok{;}
        \DecValTok{VCONST:}\NormalTok{ next\_state }\OperatorTok{=}\NormalTok{ iterm }\OperatorTok{?}\NormalTok{ DONE }\OperatorTok{:}\NormalTok{ VCONST}\OperatorTok{;}
        \DecValTok{DONE:}\NormalTok{ next\_state }\OperatorTok{=}\NormalTok{ vrchrg }\OperatorTok{?}\NormalTok{ TRLK }\OperatorTok{:}\NormalTok{DONE}\OperatorTok{;}
        \KeywordTok{default}\OperatorTok{:}\NormalTok{ next\_state }\OperatorTok{=}\NormalTok{ TRLK}\OperatorTok{;}
      \KeywordTok{endcase} \CommentTok{// case (state)}
    \KeywordTok{end}
\end{Highlighting}
\end{Shaded}

\begin{Shaded}
\begin{Highlighting}[]
   \CommentTok{//{-} Control output signals}
\NormalTok{   always\_ff }\OperatorTok{@(}\KeywordTok{posedge}\NormalTok{ clk }\DataTypeTok{or} \KeywordTok{posedge}\NormalTok{ reset}\OperatorTok{)} \KeywordTok{begin}
      \KeywordTok{if}\OperatorTok{(}\NormalTok{reset}\OperatorTok{)} \KeywordTok{begin}
\NormalTok{         state }\OperatorTok{\textless{}=}\NormalTok{ TRLK}\OperatorTok{;}
\NormalTok{         trkl }\OperatorTok{\textless{}=} \DecValTok{1}\OperatorTok{;}
\NormalTok{         fast }\OperatorTok{\textless{}=} \DecValTok{0}\OperatorTok{;}
\NormalTok{         vconst }\OperatorTok{\textless{}=} \DecValTok{0}\OperatorTok{;}
\NormalTok{         done }\OperatorTok{\textless{}=} \DecValTok{0}\OperatorTok{;}
      \KeywordTok{end}
      \KeywordTok{else} \KeywordTok{begin}
\NormalTok{         state }\OperatorTok{\textless{}=}\NormalTok{ next\_state}\OperatorTok{;}
         \KeywordTok{case} \OperatorTok{(}\NormalTok{state}\OperatorTok{)}
           \DecValTok{TRLK:} \KeywordTok{begin}
\NormalTok{              trkl }\OperatorTok{\textless{}=} \DecValTok{1}\OperatorTok{;}
\NormalTok{              fast }\OperatorTok{\textless{}=} \DecValTok{0}\OperatorTok{;}
\NormalTok{              vconst }\OperatorTok{\textless{}=} \DecValTok{0}\OperatorTok{;}
\NormalTok{              done }\OperatorTok{\textless{}=} \DecValTok{0}\OperatorTok{;}
           \KeywordTok{end}
           \DecValTok{FAST:} \KeywordTok{begin}
\NormalTok{              trkl }\OperatorTok{\textless{}=} \DecValTok{0}\OperatorTok{;}
\NormalTok{              fast }\OperatorTok{\textless{}=} \DecValTok{1}\OperatorTok{;}
\NormalTok{              vconst }\OperatorTok{\textless{}=} \DecValTok{0}\OperatorTok{;}
\NormalTok{              done }\OperatorTok{\textless{}=} \DecValTok{0}\OperatorTok{;}

           \KeywordTok{end}
           \DecValTok{VCONST:} \KeywordTok{begin}
\NormalTok{              trkl }\OperatorTok{\textless{}=} \DecValTok{0}\OperatorTok{;}
\NormalTok{              fast }\OperatorTok{\textless{}=} \DecValTok{0}\OperatorTok{;}
\NormalTok{              vconst }\OperatorTok{\textless{}=} \DecValTok{1}\OperatorTok{;}
\NormalTok{              done }\OperatorTok{\textless{}=} \DecValTok{0}\OperatorTok{;}

           \KeywordTok{end}
           \DecValTok{DONE:} \KeywordTok{begin}
\NormalTok{              trkl }\OperatorTok{\textless{}=} \DecValTok{0}\OperatorTok{;}
\NormalTok{              fast }\OperatorTok{\textless{}=} \DecValTok{0}\OperatorTok{;}
\NormalTok{              vconst }\OperatorTok{\textless{}=} \DecValTok{0}\OperatorTok{;}
\NormalTok{              done }\OperatorTok{\textless{}=} \DecValTok{1}\OperatorTok{;}
           \KeywordTok{end}
         \KeywordTok{endcase} \CommentTok{// case (state)}
      \KeywordTok{end} \CommentTok{// else: !if(reset)}
   \KeywordTok{end}
\KeywordTok{endmodule}
\end{Highlighting}
\end{Shaded}

\subsubsection{Synthesize FSM with
yosys}\label{synthesize-fsm-with-yosys}

\emph{dicex/sim/verilog/bcharger\_sv/bcharger.ys}

\begin{Shaded}
\begin{Highlighting}[]

\CommentTok{\# read design}
\NormalTok{read\_verilog}\OtherTok{ {-}sv}\NormalTok{ bcharger.sv;}
\NormalTok{hierarchy}\OtherTok{ {-}top}\NormalTok{ bcharger;}

\CommentTok{\# the high{-}level stuff}
\NormalTok{fsm; opt; memory; opt;}

\CommentTok{\# mapping to internal cell library}
\NormalTok{techmap; opt;}
\NormalTok{synth;}
\NormalTok{opt\_clean;}

\CommentTok{\# mapping flip{-}flops }
\NormalTok{dfflibmap  {-}liberty ../../../lib/SUN\_TR\_GF130N.lib}

\CommentTok{\# mapping logic }
\NormalTok{abc}\OtherTok{ {-}liberty}\NormalTok{ ../../../lib/SUN\_TR\_GF130N.lib}

\CommentTok{\# write synth netlist}
\NormalTok{write\_verilog bcharger\_netlist.v}
\NormalTok{read\_verilog  ../../../lib/SUN\_TR\_GF130N\_empty.v}
\NormalTok{write\_spice}\OtherTok{ {-}big\_endian {-}neg}\NormalTok{ AVSS}\OtherTok{ {-}pos}\NormalTok{ AVDD}\OtherTok{ {-}top}\NormalTok{ bcharger bcharger\_netlist.sp}

\CommentTok{\# write dot so we can make image}
\OtherTok{show {-}format}\NormalTok{ dot}\OtherTok{ {-}prefix}\NormalTok{ bcharger\_synth}\OtherTok{ {-}colors} \DecValTok{1}\OtherTok{ {-}width {-}stretch}
\NormalTok{clean}
\end{Highlighting}
\end{Shaded}


{
\centering
\aicfig{media/bcharger_synth.pdf}
}

\emph{Figure 71: Gate level netlist of the battery charger after yosys
synthesis to the SUN TR GF130N cell library, with six flip-flops and the
next-state logic between them}

\section{Summary}\label{summary}

The one-page version of this chapter:

\begin{itemize}
\tightlist
\item
  CMOS logic is a PMOS pull-up network against an NMOS pull-down: the
  inverter is the atom
\item
  Delay is RC: Elmore for a hand estimate, logical effort for sizing -
  stage long paths at an effort around four
\item
  Registers, combinational clouds and a clock make synchronous design;
  setup and hold are checked at every capture
\item
  Static timing analysis walks every path over PVT: positive slack, or
  it does not ship
\item
  Dynamic power is CV\^{}2 f; leakage is what you pay even when nothing
  switches
\end{itemize}

\section{Would you like to know
more?}\label{would-you-like-to-know-more}

Weste and Harris, \emph{CMOS VLSI Design} - the standard treatment of
logic, timing and power

Sutherland, Sproull and Harris, \emph{Logical Effort} - the short book
that made stage sizing a method rather than a habit

\phantomsection\label{refs}
\begin{CSLReferences}{0}{0}
\bibitem[\citeproctext]{ref-wulff17}
\CSLLeftMargin{{[}1{]} }%
\CSLRightInline{C. Wulff and T. Ytterdal, {``A compiled 9-bit 20-MS/s
3.5-fJ/conv.step SAR ADC in 28-nm FDSOI for bluetooth low energy
receivers,''} \emph{IEEE Journal of Solid-State Circuits}, vol. 52, no.
7, pp. 1915--1926, 2017, doi:
\href{https://doi.org/10.1109/JSSC.2017.2685463}{10.1109/JSSC.2017.2685463}.
}

\end{CSLReferences}
